\subsubsection{Hardwares Introduction}

Modern extended reality (XR) hardware can be classified based on the location of its compute unit into two categories: PC-tethered virtual reality (PCVR) systems and standalone (all-in-one) systems.

PCVR (Tethered): The headset acts as a peripheral. Assets are rendered on a high-end PC and streamed to the device via a wired (Link) or wireless (Air Link) connection. This allows developers to utilize the full power of engines like Unreal Engine, enabling high-fidelity features such as Lumen and Nanite. PC VR also supports specialized peripherals like haptic suits or omnidirectional treadmills, as seen in industry-standard titles like Half-Life: Alyx.

Standalone (All-in-One): Standalone headsets are self-contained systems with internal mobile chipsets. These provide high mobility but require significant optimization. High-end features must often be disabled or simplified to maintain performance. However, lightweight engines like Unity or Godot excel in this space, offering a balance between visual quality and mobile performance.

\parencite{rendevski2022vr} note that Standalone VR (SVR) applications offer significant advantages over PC VR. With proper optimization techniques, almost any PC VR scenario can be developed and deployed as SVR while still providing an acceptable user experience. Differently, a recent survey from \parencite{gdc2026report} shows that among developers interested in creating VR/AR/MR games, the SteamVR platform (PC VR) slightly surpassed the Meta Quest / Horizon Store (68\% vs. 64\%) as their platform of interest. These were followed by PlayStation VR/VR2 (34\%) and Apple visionOS (15\%). However, when filtering for developers who have actually worked on VR/AR/MR games, a strong preference remains for the Meta Quest (72\%) over SteamVR (59\%).

Due to hardware availability, this paper focuses exclusively on development for the Meta Quest 3.For Head-Mounted Devices (HMDs) like the Meta Quest, two primary architectural configurations exist:
